
% ════════════════════════════════════════════
% 具体分析（5.3.1.具体分析.tex，三级子块）
% ════════════════════════════════════════════
\subsubsection{滚动调整的动因与实施思路}

本问属于滚动预测下的随机调度与调整优化问题，采用滚动时域线性规划方法，先将四时刻发布的小时光伏预报插值到 10 分钟分辨率，再在 0:00、6:00、12:00 与 18:00 四个预报时刻依次制定计划购电与调整购电策略，并引入违约惩罚与超额电价机制。决策过程分三个层次：在 0:00 依据最新预报制定全天计划购电策略；在 6:00、12:00、18:00 依据不断更新的预报对剩余时段进行滚动调整，其中计划购电量高于调整购电量的部分按交易电价 50\% 计违约费用，调整购电量高于计划购电量的部分按交易电价 1.5 倍计价；再以实际供电缺口核算 5 倍电价的紧急购电费用，得到全年总购电费用并分析引入其他预报时刻调整的收益（1 月为预热期，费用汇总仅统计 2 月 1 日—12 月 31 日）。
